% !PRINTED 546
\weberpage{546}
%<<CONT
\[
\begin{aligned}
m &= 14,\ 34,\ 46,\ 82,\ 142 \equiv \pm2 \quad (\mathrm{mod}\ 16),\\
m &= 17,\ 49,\ 73,\ 97,\ 193 \equiv 1 \qquad (\mathrm{mod}\ 8).
\end{aligned}
\]
Die Formenklassen des \gesperrt{Hauptgeschlechts} können in diesen Fällen repräsentiert werden für ein gerades $m$ durch

\weq{1}{(1, 0, m), \quad \left(2, 0, \dfrac{m}{2}\right),}

für ein ungerades $m$ durch

\weq{2}{(1, 0, m), \quad \left(2, 1, \dfrac{m+1}{2}\right),}

von denen die letztere äquivalent ist mit

\weq{3}{\left(\dfrac{m+1}{2},\ \dfrac{m-1}{2},\ \dfrac{m+1}{2}\right).}

Für die Formen (1) sind nach § 127, 6. die Klasseninvarianten
\[
\frac{f_1(\sqrt{-m})^2}{\sqrt2} \quad \text{und} \quad \frac{1}{\sqrt2} f_2\left(\frac{\sqrt{-m}}{2}\right)^2 = \frac{\sqrt2}{f_1(\sqrt{-m})^2} \qquad (\S\ 34)
\]
und für die Formen (2), (3) (§ 127, 3.)
\[
\frac{f(\sqrt{-m})^2}{\sqrt2} \quad \text{und} \quad \frac{1}{\sqrt2} f\left(\frac{1+\sqrt{-m}}{1-\sqrt{-m}}\right)^2 = \frac{\sqrt2}{f(\sqrt{-m})^2},
\]
und nach § 142 sind dies ganze algebraische Zahlen.

Setzen wir also, entsprechend den beiden Fällen:

\weq{4}{\sqrt2\, x = f_1(\sqrt{-m})^2, \quad f(\sqrt{-m})^2,}

so ist
\[
x + \frac{1}{x}
\]
eine ganze algebraische Zahl, die nur eine Quadratwurzel enthält, und aus § 138 erhalten wir Aufschluß, welche Quadratwurzel darin vorkommt.

Es ist $\sqrt2$, wenn $m \equiv 6$ $(\mathrm{mod}\ 8)$ ist, also für $m = 14, 46, 142$, $\sqrt{\dfrac{m}{2}}$, wenn $m \equiv 2$ $(\mathrm{mod}\ 8)$ ist, also für $m = 34, 82$, ferner $\sqrt m$ im Fall eines ungeraden $m$ (mit Ausnahme von $m = 49$, wo $\sqrt7$ an die Stelle tritt).

Setzen wir also

\weq{5}{x + \frac{1}{x} = a + b\sqrt p,}

so sind $a, b$ rationale Zahlen, die höchstens den Nenner 2 haben. Es müssen aber auch $a$ und $b$ positiv sein. Denn ändern wir in (5) das Vorzeichen von $\sqrt p$, so entsteht eine andere quadratische
%CONT>>
